\documentclass[11pt]{article}
\usepackage[a4paper,margin=1in]{geometry}
\usepackage{amsmath,amssymb,amsthm}
\usepackage{graphicx}
\usepackage{hyperref}

\newtheorem{theorem}{Theorem}

\title{A Minimal-Operator-Space Subcase of Local Reflexivity for Completely \(M\)-Embedded Spaces}
\author{fa\_banach\_001 lane 7}
\date{June 18, 2026}

\begin{document}
\maketitle

\paragraph{Status.}
This is a \texttt{partial\_result\_likely\_valid} for a question in
Sharma's paper \cite{Sharma2009}.  It gives an affirmative answer for the
minimal operator-space structures arising from classical \(M\)-embedded Banach
spaces.  It is not a full answer for arbitrary one-sided \(M\)-embedded
operator spaces.

\section*{Source Question}

On page 12, Sharma asks whether right \(M\)-embedded or completely
\(M\)-embedded operator spaces are locally reflexive.

\begin{center}
\includegraphics[width=\textwidth]{figures/open_problem_crop.png}
\end{center}

\section*{Proof Intuition}

The source paper already supplies many completely \(M\)-embedded examples:
if \(E\) is a classical \(M\)-embedded Banach space, then \(\operatorname{Min}(E)\)
is completely \(M\)-embedded as an operator space.  For these examples local
reflexivity comes from a different and very robust mechanism: minimal operator
spaces sit completely isometrically inside commutative \(C^*\)-algebras.
Commutative \(C^*\)-algebras are nuclear, hence exact, and exactness passes to
operator subspaces.  Since exact operator spaces are locally reflexive, the
minimal examples are locally reflexive without using any special
one-sided \(M\)-ideal structure beyond Sharma's observation that they are
indeed completely \(M\)-embedded.

\section*{Theorem}

\begin{theorem}
Let \(E\) be an \(M\)-embedded Banach space.  Then the minimal operator space
\(\operatorname{Min}(E)\) is completely \(M\)-embedded and locally reflexive.
Indeed, \(\operatorname{Min}(E)\) is \(1\)-exact.
\end{theorem}

\begin{proof}
Sharma records that if \(E\) is an \(M\)-embedded Banach space, then
\(\operatorname{Min}(E)\) is completely \(M\)-embedded as an operator space
\cite[Section 2]{Sharma2009}.  It remains only to prove local reflexivity.

Let \(K=\operatorname{Ball}(E^*)\) with its weak-star topology.  The canonical
Banach-space embedding
\[
  E \longrightarrow C(K), \qquad x\mapsto (\varphi\mapsto \varphi(x)),
\]
is isometric.  By the defining property of the minimal operator-space
structure, this identifies \(\operatorname{Min}(E)\) completely isometrically
with an operator subspace of the commutative \(C^*\)-algebra \(C(K)\).

Commutative \(C^*\)-algebras are nuclear and therefore \(1\)-exact.  Exactness
passes to operator subspaces, so \(\operatorname{Min}(E)\) is \(1\)-exact.
Finally, exact operator spaces are locally reflexive; this is the standard
exactness-to-local-reflexivity implication used in operator-space theory and
also invoked in the source remark's phrase ``\(1\)-exact, and hence locally
reflexive'' \cite[p.~12]{Sharma2009}.  Hence \(\operatorname{Min}(E)\) is
locally reflexive.
\end{proof}

\section*{Verification and Novelty Notes}

The proof is a chain of standard implications and contains no computational
component.  The verification focus is whether one accepts the standard
operator-space facts: minimal spaces embed into commutative \(C^*\)-algebras,
commutative \(C^*\)-algebras are \(1\)-exact, exactness passes to subspaces,
and exact operator spaces are locally reflexive.

The cheap indexes were searched for arXiv:0902.4257, one-sided
\(M\)-embedded spaces, completely \(M\)-embedded spaces, local reflexivity,
exactness, and \(C_\infty(X)\).
Bounded web searches for exact local-reflexivity phrases found the source
paper but no separate full resolution.  Since the proof uses standard
operator-space facts, novelty confidence is low-to-moderate; the packet is best
viewed as a useful explicit subcase and duplicate-memory entry.

\section*{Limitations}

This does not answer whether all right \(M\)-embedded or completely
\(M\)-embedded operator spaces are locally reflexive.  It only treats the
minimal operator-space structures \(\operatorname{Min}(E)\) coming from
classical \(M\)-embedded Banach spaces.  In particular, it does not address the
non-minimal and genuinely one-sided nonselfadjoint examples emphasized in
\cite{Sharma2009}.

\section*{Human Review Recommendation}

Reviewers should check that the packet's operator-space background facts are
being used in their standard completely isometric form.  The mathematical
claim is intentionally modest; the main value is explicitly connecting the
source question with the minimal examples.

\begin{thebibliography}{9}
\bibitem{Sharma2009}
S. Sharma,
\emph{Operator Spaces which are One-sided \(M\)-Ideals in their bidual},
arXiv:0902.4257, 2009.
\end{thebibliography}

\end{document}
